\chapter{Lecture 13}

\hfill {\footnotesize \emph{Date: 25.04.2026}}
\todo{This part of the lecture is yet to be updated!}
\section{Primary Decomposition---continued}
\begin{proposition}
  Let \( \mathfrak{a} \) be a decomposable ideal. Let \( \mathfrak{a} = \bigcap_{i=1}^n \mathfrak{q}_i \) be a minimal primary decomposition with \( r(\mathfrak{q}_i) = \mathfrak{p}_i \). Then
  \[
    \bigcup_{i=1}^n \mathfrak{p}_i = \{x \in R \mid (\mathfrak{a}: x) \neq \mathfrak{a}\}.
  \]
  If the zero ideal is decomposable then the set \( D \) of zero divisors of  \( R \) is the union of the prime ideal belonging to  \( (0) \).
\end{proposition}

\begin{proof}
First observe that \(x+\mathfrak a\) is a zero divisor in \(R/\mathfrak a\) precisely when there exists \(y\notin \mathfrak a\) such that \(xy\in \mathfrak a\). This is equivalent to saying that there exists \(y\notin \mathfrak a\) with \(y\in(\mathfrak a:x)\). Hence
\[
    x+\mathfrak a \text{ is a zero divisor in } R/\mathfrak a
    \iff
    (\mathfrak a:x)\neq \mathfrak a.
\]

So it is enough to prove that the elements \(x\in R\) for which \(x+\mathfrak a\) is a zero divisor are exactly the elements of \(\bigcup_i\mathfrak p_i\).

We first prove the result in the special case \(\mathfrak a=(0)\). Suppose
\[
    (0)=\bigcap_{i=1}^n \mathfrak q_i
\]
is a minimal primary decomposition, where \(r(\mathfrak q_i)=\mathfrak p_i\). Let \(D\) denote the set of zero divisors of \(R\). We claim that
\[
    D=\bigcup_{i=1}^n \mathfrak p_i.
\]

First let \(x\in \mathfrak p_i\) for some \(i\). Since \(\mathfrak p_i=r(\mathfrak q_i)\), there is some \(n > 0\) such that \(x^n\in \mathfrak q_i\).

By minimality, \(\mathfrak q_i\) is not redundant. Hence
\[
    \bigcap_{j\neq i}\mathfrak q_j \nsubseteq \mathfrak q_i.
\]
Choose \(a_i\in\bigcap_{j\neq i}\mathfrak q_j\) such that \(a_i\notin\mathfrak q_i\).

Since \(x^n\in\mathfrak q_i\), we have \(x^n a_i\in\mathfrak q_i\). Choose the smallest positive integer \(m\) such that \(x^m a_i\in\mathfrak q_i\). Then \(x^{m-1}a_i\notin\mathfrak q_i\), so in particular \(x^{m-1}a_i\neq 0\).

Now \(a_i\in\mathfrak q_j\) for every \(j\neq i\), so \(x^m a_i\in\mathfrak q_j\) for every \(j\neq i\). Also \(x^m a_i\in\mathfrak q_i\). Therefore
\[
    x^m a_i\in \bigcap_{j=1}^n \mathfrak q_j=(0).
\]
Thus \(x(x^{m-1}a_i)=0\), where \(x^{m-1}a_i\neq 0\). Hence \(x\) is a zero divisor. Therefore \(\bigcup_i\mathfrak p_i\subseteq D\).

Conversely, let \(x\in D\). Then there exists \(0\neq y\in R\) such that \(xy=0\). Since
\[
    (0)=\bigcap_{i=1}^n \mathfrak q_i,
\]
the condition \(y\neq 0\) implies that \(y\notin\mathfrak q_k\) for at least one \(k\). But \(xy=0\in\mathfrak q_k\). Since \(\mathfrak q_k\) is primary and \(y\notin\mathfrak q_k\), it follows that \(x^m\in\mathfrak q_k\) for some \(m > 0\). Therefore \(x\in r(\mathfrak q_k)=\mathfrak p_k\). Hence \(D\subseteq\bigcup_i\mathfrak p_i\).

Thus, when \(\mathfrak a=(0)\), the set of zero divisors of \(R\) is exactly \(\bigcup_i\mathfrak p_i\).

Now return to the general case. Let
\[
    \pi:R\to R/\mathfrak a, r\mapsto r+\mathfrak a
\]
be the quotient map. Since \(\mathfrak a=\bigcap_i\mathfrak q_i\), we have
\[
    (0)=\bigcap_{i=1}^n \pi(\mathfrak q_i)
\]
inside \(R/\mathfrak a\).

We claim that each \(\pi(\mathfrak q_i)\) is primary in \(R/\mathfrak a\). Suppose
\[
    (r_1+\mathfrak a)(r_2+\mathfrak a)\in \pi(\mathfrak q_i).
\]
Then \(r_1r_2+\mathfrak a\in\pi(\mathfrak q_i)\). So there exists \(q_i\in\mathfrak q_i\) such that \(r_1r_2+\mathfrak a=q_i+\mathfrak a\). Hence \(r_1r_2-q_i\in\mathfrak a\subseteq\mathfrak q_i\). Since \(q_i\in\mathfrak q_i\), we get \(r_1r_2\in\mathfrak q_i\).

Because \(\mathfrak q_i\) is primary, \(r_1\in\mathfrak q_i\), or \(r_2^n\in\mathfrak q_i\) for some \(n > 0\). Therefore
\[
  r_1+\mathfrak a\in\pi(\mathfrak q_i) \text{ or }
    (r_2+\mathfrak a)^N\in\pi(\mathfrak q_i).
\]
Thus \(\pi(\mathfrak q_i)\) is primary in \(R/\mathfrak a\).

Moreover,
\[
    r(\pi(\mathfrak q_i))
    =
    \pi(r(\mathfrak q_i))
    =
    \pi(\mathfrak p_i).
\]
\todo{To be written...}
% The decomposition remains minimal in the quotient. Indeed, the radicals \(\pi(\mathfrak p_i)\) are distinct because the \(\mathfrak p_i\)'s are distinct and all contain \(\mathfrak a\). Also, if \(\pi(\mathfrak q_i)\) were redundant, then taking inverse images under \(\pi\) would imply that \(\mathfrak q_i\) was redundant in the original decomposition.
%
% So
% \[
%     (0)=\bigcap_{i=1}^n \pi(\mathfrak q_i)
% \]
% is a minimal primary decomposition of the zero ideal in \(R/\mathfrak a\). Applying the special case to the ring \(R/\mathfrak a\), the set of zero divisors of \(R/\mathfrak a\) is
% \[
%     \bigcup_{i=1}^n \pi(\mathfrak p_i).
% \]
%
% Thus \(x+\mathfrak a\) is a zero divisor in \(R/\mathfrak a\) if and only if \(x\in\mathfrak p_i\) for some \(i\). Therefore the set of zero divisors of \(R/\mathfrak a\), viewed inside \(R\), is
% \[
%     \bigcup_{i=1}^n \mathfrak p_i.
% \]
% Equivalently,
% \[
%     \bigcup_{i=1}^n \mathfrak p_i
%     =
%     \{x\in R:(\mathfrak a:x)\neq\mathfrak a\}.
% \]
% This completes the proof.
\end{proof}

\begin{remark}
A useful fact is the following. If \(f:R\to R'\) is a ring homomorphism and \(\mathfrak q'\) is a \(\mathfrak p'\)-primary ideal of \(R'\), then \(f^{-1}(\mathfrak q')\) is \(f^{-1}(\mathfrak p')\)-primary in \(R\).

Indeed, if \(xy\in f^{-1}(\mathfrak q')\), then \(f(x)f(y)\in\mathfrak q'\). Since \(\mathfrak q'\) is primary, either \(f(x)\in\mathfrak q'\), or \(f(y)^N\in\mathfrak q'\) for some \(N\geq 1\). Hence either \(x\in f^{-1}(\mathfrak q')\), or \(y^N\in f^{-1}(\mathfrak q')\). Therefore \(f^{-1}(\mathfrak q')\) is primary, and
\[
    r(f^{-1}(\mathfrak q'))
    =
    f^{-1}(r(\mathfrak q'))
    =
    f^{-1}(\mathfrak p').
\]
\end{remark}

